% Generated by scripts/build-paper.py — edit paper/src/survey.src.md
\section{Introduction}\label{sec:introduction}

Many software decisions are small and bounded: route a ticket, accept or escalate a judgment, pick the next tool, code a variable from a narrative, admit a network policy. Large language models can make such decisions, but they return strings that must be parsed, validated and retried, and the confidence they state is usually a verbalised estimate rather than a probability a program can threshold \citep{arxiv2207_05221,arxiv2503_15850}. Typed decision models change the contract. The caller declares the admissible answers; the model returns one of them with a probability for every option; no free text is generated, so a parse failure is impossible by construction.

TypeSafe AI's Jev, released in early access on 15 September 2026, made this contract prominent \citep{typesafe2026launch}. Its launch materials describe a new class of models trained with “Reinforcement Learning for Calibrated Decisions” (RLCD) and report end-to-end responses of 70–500 ms at \$0.042 per million input tokens \citep{typesafe2026launch,typesafe2026models}. Three weeks later the literature already spanned audits of option semantics and calibration, open models trained to copy the interface, and systems in judging, annotation, moderation, security, networks, databases, search, robotics, medicine and materials discovery. The claims range from schema validity to probability coherence, cost, latency and end-task success; the protocols range from 20 curated scientific choices to half a million police crash narratives.

Others have begun to map this literature. An early evidence audit reviewed the 28 papers posted between 19 and 24 September and derived a 14-item evaluation checklist \citep{arxiv2609_32160}; a systematisation of knowledge examined semantic decision engines in network control loops \citep{arxiv2610_06425}; a public survey draft places typed decisions in the history of classification, calibration and routing \citep{decisionsnottokens2026}; and community lists catalogue models, projects and benchmarks \citep{gh_omnijev_awesome_jev_gallery,gh_omnijev_awesome_jev_papers}. This survey is an \textbf{empirical survey and evidence audit}: it reads every core study in full, re-reads studies when they post new versions, records each checkable claim with its locator and scope, and keeps incomparable numbers apart.

We keep three objects distinct throughout: (A) the \textbf{commercial model}, TypeSafe's hosted Jev, whose architecture and training data are undisclosed; (B) the \textbf{paradigm}, Jev-like typed decision models and readouts built by others; and (C) \textbf{systems} whose results depend on typed decisions. A result about one is not evidence about another.

We ask six questions. \textbf{RQ1 (definition and lineage):} how do runtime-defined finite-option decisions relate to classifiers, zero-shot labelling, rerankers, reward models and constrained generation? \textbf{RQ2 (probability quality):} what do returned probabilities and confidence mean, where do they hold, and are they coherent across related questions? \textbf{RQ3 (efficiency attribution):} how much of a speed or cost gain comes from the readout, and how much from model size, batching, prefix sharing, caching, pruning or the network? \textbf{RQ4 (selective control):} which decisions can be delegated, and when should software escalate, abstain or let code compute? \textbf{RQ5 (robustness and security):} how do option names, order, definitions, missing options, language, injected text and poisoned checkpoints change decisions? \textbf{RQ6 (openness):} what do alternatives release, and how much of the literature survives a common protocol?

Our contributions are concrete artefacts rather than claims of priority:

\begin{enumerate}[leftmargin=*,itemsep=2pt]
  \item A versioned evidence base of 175 references and 215 evidence records. Each record carries a source locator, evidence type, test level, measurement scope, sample size and model version, with explicit reasons where a value is missing (§3).
  \item A synthesis into ten findings, each listing supporting \emph{and} qualifying evidence (§7). Study families are synthesised once, and nothing is pooled across incompatible scopes.
  \item An implementation lineage of six readout families and an audit of the open ecosystem that separates code, weights, training code, evaluation material, data and raw predictions (§6, §9).
  \item A catalogue of failure modes with sources and mitigations (§8), and a proposed evaluation protocol whose first experiments target binding, rejection, coherence and escalation together (§10).
  \item A public website, generated README, BibTeX and CSV exports, all produced from the same data by scripts with validation and link checks.
\end{enumerate}
